\section{TD3 --- Compression itérative}\label{sec:td3-compression}

\begin{td-exo}[Profondeur bornée ne signifie pas \algo{FPT}]
    Deux procédures ont une profondeur au plus \(k\).
    La première crée au plus deux fils par nœud, la seconde au plus \(n\).
    Le travail local coute \(n^2\).

    \begin{enumerate}
        \item Borner leurs temps de calcul et classer ces bornes comme \algo{FPT} ou \algo{XP}.

        \item Pour une constante \(c\) fixée et \(k > c\), montrer qu'aucune constante \(f(k)\) ne majore \(n^{k-c}\) pour tous les entiers \(n\).

        \item Un algorithme en temps \(n^k\) prouve-t-il que le problème n'est pas \algo{FPT} ?

        \item Si le nombre de fils est \(k + 1\), la conclusion change-t-elle ?
    \end{enumerate}
\end{td-exo}

\begin{td-sol}
    On note \(p_1\) la procédure qui crée au plus deux fils par nœud et \(p_2\) celle qui en crée au plus \(n\).

    \begin{enumerate}
        \item Le nombre de feuilles pour \(p_1\) est au plus
        \begin{equation*}
            {b_1}^{d_1}
        \end{equation*}
        où \(b_1\) est le facteur de branchement et \(d_1\) la profondeur.

        Ici, on a \(b_1 \leq 2\) et \(d_1 \leq k\) donc le nombre de feuilles \(f_1\) respecte 
        \begin{equation*}
            \begin{aligned}
                f_1 
                &\leq {b_1}^{d_1} \\
                &\leq 2^k
            \end{aligned}
        \end{equation*}
        Ainsi, le nombre de nœuds dans l'arbre de recherche pour la procédure \(p_1\) est au plus \(2f = 2^{k+1}\).
        Le temps de calcul en chaque nœud étant exactement \(n^2\), la procédure \(p_1\) finit en temps au plus \(2^{k+1}\cdot n^2\).

        En posant \(f(k) = 2^{k+1}\), on obtient que \(p_1\) termine en \(f(k) n^2\) donc la procédure \(p_1\) est \algo{FPT}.
        
        \vspace{\baselineskip}

        On procède de même pour \(p_2\) :
        \begin{equation*}
            {b_2}^{d_2}
        \end{equation*}
        avec \(b_2 = n\) et \(d_2 = k\).

        On a alors un temps de calcul en au plus \(2n^k \cdot n^2 = 2 n^{k+2}\). 
        Le facteur \(k\) étant dans l'exposant de \(n\), la procédure \(p_2\) est \algo{XP}.

        \item À \(k\) fixé, \(f(k)\) est une constante. Or \(n^{k-c}\) est une fonction avec \(\displaystyle\lim_{n\to\infty} n^{k-c} = +\infty\).
        Donc aucune constante ne majore cette fonction pour tout \(n\).

        \item Un tel algorithme ne suffit pas à prouver qu'un problème n'est pas \algo{FPT}, pour cela, il faut prouver qu'il est \algo{W}\([1]\) difficile.

        \item Si le nombre de fils à chaque nœud devient \(k+1\), le nombre de feuilles devient \(2^{k+1}\) qui est toujours \(\BigO{2^k}\).
        La procédure est donc toujours \algo{FPT}.
    \end{enumerate}
\end{td-sol}

\begin{td-exo}[Closest String]
    On reprend les mots \(s_1,\ldots,s_r \in \Sigma^L\), le rayon \(d\) et le pseudo-code du cours.
    La distance de Hamming \(d_H(x, y)\) est le nombre de positions où \(x\) et \(y\) diffèrent.
    L'appel initial se fait avec \((z, l) =  (s_1, d)\).

    \begin{enumerate}
        \item Soit \(c\) une solution et \(s_i\) un mot avec \(d_H(z, s_i) > d\).
        Dans un ensemble \(P\) de \(d+1\) positions de désaccord, montrer qu'il existe \(p \in P\) tel que \(c[p] = s_i[p] \neq z[p]\).
        Pourquoi ne suffit-il pas de dire seulement que \(c\) et \(z\) diffèrent quelque part dans \(P\) ?

        \item Montrer que, si \(d_H(z, c) \leq l\), au moins un appel enfant vérifie \(d_H(z^{(p)}, c) \leq l-1\).
        Justifier l'invariant à la racine.

        \item En déduire que l'algorithme trouve une solution si elle existe et n'en renvoie jamais une fausse.
        Traiter le cas \(l = 0\).

        \item Pour \(d\geq 1\), borner le nombre de nœuds et obtenir un temps \(\BigO{(d+1)^d r L}\).
        Traiter séparément \(d = 0\).
    \end{enumerate}
\end{td-exo}